% !TeX root = ./chapter_1_main.tex
\section{Optical Fibers}
A waveguide is a device that uses dielectric media with high refractive index to confine light within a region of space inhibiting its natural spread typical of free space propagation. The wave that propagates inside a waveguide is called guided mode (or guided wave).

Optical fibers are circularly symmetric waveguide, and are by far the most common waveguide in the world. Following chapter III of \citeauthor*{yariv2007photonics} \cite{yariv2007photonics}, this chapter is devoted to understanding the physical principles that allow optical fibers to guide confined light, with particular focus on the shape of their modes of propagation.

%* ##########################################################
\subsection{Circularly symmetric wave guides}
%* ##########################################################
An optical fiber is model as a cylinder called core with refractive index $n_1$ surrounded by a hollow cylinder called cladding of refractive index $n_2$. The refractive index has therefore the step profile
\begin{equation}
    n(r)=
    \begin{cases}
        n_1, & 0 < r < a \\
        n_2, & a < r < b,
    \end{cases}
\end{equation}
where $a$ is the radius of the core and $b$ is the outer radius of the cladding. Usually the radius of the cladding is much larger that the radius of the core ($b\gg a$), so that the electric field can be considered zero at $r=b$. For example in fiber that guide a single mode of light with $\lambda=1.5~\mu\mathrm{m}$, the typical core radius is $a=5~\mu\mathrm{m}$ while the cladding external radius is in the order of $b=100~\mu\mathrm{m}$.

\begin{figure}[h]
    \centering
    \import{./tikz}{optical_fiber_scheme.tex}
    \caption{Schematic representation of a step-index optical fiber. (left) and wave vector diagram showing the geometric relationship between the wave number in the core $n_1 k_0$, the longitudinal propagation constant $\beta$, and the transverse wave number $h$ (right).}
    \label{fig:optical_fiber_scheme}
\end{figure}

\subsubsection{Cylindrical symmetric wave equation}
Given the symmetry of the system, it is natural to work with the cylindrical coordinates:
\begin{equation}
    \begin{cases}
        \mathbf{u_r}    = & \hat{\mathbf{x}}\cos\phi  + \hat{\mathbf{y}}\sin\phi  \\
        \mathbf{u_\phi} = & - \hat{\mathbf{x}}\sin\phi + \hat{\mathbf{y}}\cos\phi \\
        \mathbf{u_z}    = & \hat{\mathbf{z}},
    \end{cases}
\end{equation}
thus the electric and magnetic fields are expressed in terms of the components $E_r$, $E_\phi$, $E_z$, $H_r$, $H_\phi$ and $H_z$. In cylindrical coordinates the unit vectors $\mathbf{u_r}$ and $\mathbf{u_\phi}$ are not constant vectors ($\partial\mathbf{u_r}/\partial \phi = \mathbf{u_\phi}$ and $\partial \mathbf{u_\phi} /\partial \phi = -\mathbf{u_r}$), this complicates the wave equation for the associated components of the field. Nevertheless, as $\mathbf{u_z}=\mathbf{\hat{z}}$ the equation for the z components remain simple:
\begin{equation}
    \label{eq:z_component_wave_equation}
    (\nabla^2 + k^2)
    \begin{bmatrix}
        E_z \\
        H_z
    \end{bmatrix}
    =0,
\end{equation}
where $k^2=\omega^2n/c^2$ and $\nabla^2$ is the Laplacian operator in cylindrical coordinates.
To solve the cylindrical wave equation, the usual approach is to solve for $E_z$ and $H_z$ and then express $E_r$, $E_\phi$, $H_r$ and $H_\phi$ in terms of $E_z$ and $H_r$, this is possible thanks to Maxwell's curl equations in cylindrical coordinates that allows the $r$ and $\phi$ components to be expressed as a function of solely $E_z$ and $H_z$ (Appendix A of \cite{yariv2007photonics}).

Assuming for $E_z$ and $H_z$ a time $t$ dependence of $e^{i(\omega t-\beta z)}$
\begin{equation}
    \label{eq:z_field}
    \begin{bmatrix}
        \mathbf{E}(\mathbf{r},t) \\
        \mathbf{H}(\mathbf{r},t)
    \end{bmatrix}
    =
    \begin{bmatrix}
        \mathbf{E}(r, \phi) \\
        \mathbf{H}(r, \phi)
    \end{bmatrix}
    e^{i(\omega t -\beta z)},
\end{equation}
using the cylindrical Laplacian, equation \eqref{eq:z_component_wave_equation}, one gets:
\begin{equation}
    \label{eq:z_component_effective_equation}
    \left(
    \frac{\partial^2}{\partial r^2} +
    \frac{1}{r} \frac{\partial^2}{\partial r} +
    \frac{1}{r^2}\frac{\partial^2}{\partial\phi^2}+
    (k^2 - \beta^2)
    \right)
    \begin{bmatrix}
        E_z \\
        H_z
    \end{bmatrix}
    =0.
\end{equation}

Equation \eqref{eq:z_component_effective_equation} is solved by assuming the separation of variable
\begin{equation}
    F_z(r, \phi)=R(r)\Phi(\phi),
\end{equation}
where $F_z=E_z, H_z$, obtaining
\begin{equation}
    \Phi \frac{d^2R}{dr^2} + \frac{1}{r}\Phi\frac{dR}{dr} + \frac{1}{r^2}R\frac{d^2\Phi}{d\phi^2} +(k^2-\beta^2)R(r)\Phi(\phi) =0,
\end{equation}
which is a separable differential equation.

The variables are separated by dividing by $R\Phi$, obtaining
\begin{equation}
    r^2
    \left(
    \frac{1}{R}\frac{d^2R}{dr^2}+
    \frac{1}{r}\frac{1}{R}\frac{dR}{dr}+
    (k^2 - \beta^2)
    \right)
    =
    -\frac{1}{\Phi}\frac{d^2\Phi}{d\phi^2}.
\end{equation}

The last equation is in the form $f(f)=g(\phi)$, therefore the equality holds $\forall r, \phi$ if and only if the two sides are equal to same constant value:
\begin{equation}
    \begin{cases}
        r^2
        \left(
        \frac{1}{R}\frac{d^2R}{dr^2}+
        \frac{1}{r}\frac{1}{R}\frac{dR}{dr}+
        (k^2-\beta^2)
        \right)
        = l^2 \\[3mm]
        -\frac{1}{\Phi}\frac{d^2\Phi}{d\phi^2} = l^2
    \end{cases}
\end{equation}

The angular part has a simple exponential solution:
\begin{equation}
    \Phi(\phi) = C_1e^{+il\phi} + C_2 e^{-il\phi}
\end{equation}
and since $\Phi(\phi)$ must be single valued ($\Phi(0)=\Phi(2n\pi)$) the value of $l$ is constrained to
\begin{equation}
    l = 0,1,2,...
\end{equation}

The solution can be therefore expressed as:
\begin{equation*}
    \begin{bmatrix}
        E_z \\
        H_z
    \end{bmatrix}
    =R(r)e^{\pm il\phi}
\end{equation*}
where the differential equation for the radial part is
\begin{equation}
    \label{eq:bessel_equation}
    \frac{\partial^2 R}{\partial r^2} +
    \frac{1}{r}\frac{\partial R}{\partial r} +
    \left(k^2-\beta^2-\frac{l^2}{r^2}\right) R =0
\end{equation}
which is the Bessel differential equation, whose solution is given by a combination of the Bessel functions.

\subsubsection{The Bessel functions}
Equation \eqref{eq:bessel_equation} admits different solution depending on the sign of $k^2-\beta^2$:
\begin{itemize}
    \item if $k^2 -\beta^2 >0$ the solutions are the Bessel functions of the first and second kind:
    $$ \psi(r) = C_1 J_l(hr) + C_2 Y_l(hr);$$

    \item if $k^2 -\beta^2 <0$ the solution are the modified Bessel functions of the first and second kind:
    $$ \psi(r) = C_1 I_l(qr) + C_2 K_l(qr)$$
    where $h^2 = k^2 -\beta^2$ and $q^2=\beta^2-k^2$.
\end{itemize}

\begin{figure}[hb]
    \begin{subfigure}[t]{0.49\linewidth}
        \includegraphics[width=\linewidth]{Cap_1_figures/Bessel/J_n_Bessel.pdf}
    \end{subfigure}
    \hfill
    \begin{subfigure}[t]{0.49\linewidth}
            \includegraphics[width=\linewidth]{Cap_1_figures/Bessel/Y_n_Bessel.pdf}
    \end{subfigure}
    %=========================================
    \begin{subfigure}[b]{0.49\linewidth}
        \includegraphics[width=\linewidth]{Cap_1_figures/Bessel/I_n_Modified_Bessel.pdf}
    \end{subfigure}
    \hfill
    \begin{subfigure}[b]{0.49\linewidth}
        \includegraphics[width=\linewidth]{Cap_1_figures/Bessel/K_n_Modified_Bessel.pdf}
    \end{subfigure}
    \caption{Behavior of the ordinary and modified Bessel functions for the first four integer orders, $n \in \{0, 1, 2, 3\}$. Top panels: Ordinary Bessel functions of the first kind, $J_n(x)$ (left), and second kind, $Y_n(x)$ (right). Bottom panels: Modified Bessel functions of the first kind, $I_n(x)$ (left), and second kind, $K_n(x)$ (right).}
\end{figure}

%* ##########################################################
\subsection{Weak waveguide approximation}
%* ##########################################################

%* ##########################################################
\subsection{The LP modes}
%* ##########################################################